\documentclass{article}
\usepackage{amsmath,amssymb,amsthm}
\usepackage{algorithm}
\usepackage[noend]{algpseudocode}
\usepackage{bm}
\usepackage{geometry}
\geometry{margin=1in}
\newtheorem{theorem}{Theorem}
\newtheorem{lemma}{Lemma}
\newcommand{\EE}{\mathbb{E}}
\newcommand{\RR}{\mathbb{R}}
\newcommand{\norm}[1]{\left\lVert#1\right\rVert}
\newcommand{\proj}[1]{\Pi_{#1}}
\begin{document}

\section*{Algorithm Name}
\textbf{Differentially Private Stochastic Gradient Descent (DP‑SGD) with Gaussian Noise}

\section*{Mathematical Formulation}
Let $f:\RR^{d}\to\RR$ be the empirical risk
\[
f(x)=\frac{1}{n}\sum_{i=1}^{n}\ell(x;z_i),
\]
where $\ell(\cdot;z_i)$ is convex and $L$–smooth for every data point $z_i$.
At iteration $t$ a minibatch $\mathcal{B}_t\subset\{1,\dots,n\}$ of size $b$ is drawn uniformly
(at random, without replacement is not required).  

\begin{enumerate}
\item \textbf{Clipping:} For each $i\in\mathcal{B}_t$ compute the stochastic gradient
\[
g_i^{(t)}=\nabla\ell(x_t;z_i),\qquad 
\tilde g_i^{(t)}=\frac{g_i^{(t)}}{\max\{1,\; \|g_i^{(t)}\|/C\}},
\]
so that $\|\tilde g_i^{(t)}\|\le C$.
\item \textbf{Averaging:}
\[
\bar g_t=\frac{1}{b}\sum_{i\in\mathcal{B}_t}\tilde g_i^{(t)} .
\]
\item \textbf{Noise addition:} Sample $\xi_t\sim\mathcal N(0,\sigma^{2}I_d)$ and set
\[
\hat g_t=\bar g_t+\xi_t .
\]
\item \textbf{Update:}
\[
x_{t+1}=x_t-\eta_t \hat g_t .
\]
\item \textbf{Privacy accounting:} Each iteration corresponds to a Gaussian
mechanism with sensitivity $C/b$; the overall $(\varepsilon,\delta)$‑budget is
tracked by the moments accountant (or Rényi DP) using the sequence
$\{(\eta_t,\sigma)\}_{t=0}^{T-1}$.
\end{enumerate}

\section*{Pseudocode}
\begin{algorithm}[H]
\caption{DP‑SGD with Gaussian Noise}
\begin{algorithmic}[1]
\State \textbf{Input:} data $\{z_i\}_{i=1}^{n}$, loss $\ell$, clip bound $C>0$, 
learning rates $\{\eta_t\}_{t=0}^{T-1}$, noise std $\sigma$, batch size $b$, 
initial point $x_0$.
\For{$t=0$ \textbf{to} $T-1$}
    \State Sample a minibatch $\mathcal B_t$ of size $b$ uniformly at random.
    \State $\displaystyle g_i^{(t)}\gets\nabla\ell(x_t;z_i),\;\forall i\in\mathcal B_t$.
    \State $\displaystyle \tilde g_i^{(t)}\gets 
        \frac{g_i^{(t)}}{\max\{1,\|g_i^{(t)}\|/C\}}$   \Comment{clip}
    \State $\displaystyle \bar g_t\gets\frac{1}{b}\sum_{i\in\mathcal B_t}\tilde g_i^{(t)}$.
    \State Sample $\xi_t\sim\mathcal N(0,\sigma^{2}I_d)$.
    \State $\displaystyle \hat g_t\gets\bar g_t+\xi_t$.
    \State $x_{t+1}\gets x_t-\eta_t\hat g_t$.
    \State \textbf{(Optional)} Update privacy accountant with $(C/b,\sigma)$.
\EndFor
\State \textbf{Return} $x_T$ (or the average $\bar x_T=\frac1T\sum_{t=0}^{T-1}x_t$).
\end{algorithmic}
\end{algorithm}

\section*{Dry Run Example}
Consider the one‑dimensional quadratic $f(w)=(w-3)^2$, thus $\nabla f(w)=2(w-3)$.
Take $C=10$ (no clipping), $\eta=0.1$, $\sigma=0$, batch size $b=1$, and start
with $w_0=0$.

\begin{center}
\begin{tabular}{c|c|c|c}
$t$ & $g_t=\nabla f(w_t)$ & $w_{t+1}=w_t-\eta g_t$ & $f(w_{t+1})$\\ \hline
0 & $2(0-3)=-6$ & $0-0.1(-6)=0.6$ & $(0.6-3)^2=5.76$\\
1 & $2(0.6-3)=-4.8$ & $0.6-0.1(-4.8)=1.08$ & $(1.08-3)^2=3.6864$\\
2 & $2(1.08-3)=-3.84$ & $1.08-0.1(-3.84)=1.464$ & $(1.464-3)^2=2.3650$\\
\end{tabular}
\end{center}
The iterates move toward the optimum $w^{*}=3$, and the objective value
decreases at each step.

\newpage
\section*{Assumptions}
\begin{enumerate}
\item \textbf{Convexity and Smoothness:} $f$ is convex and $L$‑smooth,
\[
f(y)\le f(x)+\langle\nabla f(x),y-x\rangle+\frac{L}{2}\|y-x\|^{2},
\quad\forall x,y\in\RR^{d}.
\]
\item \textbf{Gradient Clipping Bound:} After clipping, $\|\tilde g_i^{(t)}\|\le C$
for all $i,t$.
\item \textbf{Noise Distribution:} $\xi_t\sim\mathcal N(0,\sigma^{2}I_d)$, independent
of the data and of previous iterates. Hence $\EE[\xi_t]=0$ and
$\EE[\|\xi_t\|^{2}]=d\sigma^{2}$.
\item \textbf{Learning‑rate Schedule:} $\eta_t=\eta$ (constant) or any
non‑increasing sequence satisfying $\eta_t\le \frac{1}{L}$.
\item \textbf{Bounded Domain (optional):} There exists $D>0$ such that
$\|x_0-x^{*}\|\le D$, where $x^{*}\in\arg\min f$.
\end{enumerate}

\section*{Main Convergence Theorem}
\begin{theorem}
\label{thm:conv}
Assume the conditions in Section 4. Let $\bar x_T=\frac{1}{T}\sum_{t=0}^{T-1}x_t$
be the average iterate. For a constant stepsize $\eta\le \frac{1}{L}$,
\[
\EE\!\left[f(\bar x_T)-f(x^{*})\right]
\le
\frac{\|x_0-x^{*}\|^{2}}{2\eta T}
+\frac{\eta}{2}\bigl(C^{2}+d\sigma^{2}\bigr).
\]
In particular, choosing $\eta=\frac{D}{\sqrt{(C^{2}+d\sigma^{2})\,T}}$ yields
\[
\EE\!\left[f(\bar x_T)-f(x^{*})\right]
=O\!\left(\frac{D\sqrt{C^{2}+d\sigma^{2}}}{\sqrt{T}}\right).
\]
\end{theorem}

\section*{Theoretical Properties}
\begin{itemize}
\item \textbf{Stability:} The bound depends linearly on the variance term
$d\sigma^{2}$; larger noise (stronger privacy) degrades the rate but never
destroys convergence because the additive term is multiplied by $\eta$.
\item \textbf{Hyper‑parameter Sensitivity:} The optimal constant stepsize is
$\eta^{\star}=D/\sqrt{(C^{2}+d\sigma^{2})T}$, balancing the bias term
$\|x_0-x^{*}\|^{2}/(2\eta T)$ and the variance term $\eta(C^{2}+d\sigma^{2})/2$.
\item \textbf{Comparison to Non‑private SGD:} Setting $\sigma=0$ recovers the
standard SGD bound $\frac{D^{2}}{2\eta T}+\frac{\eta C^{2}}{2}$.
\item \textbf{Privacy–Utility Trade‑off:} The noise variance $\sigma^{2}$
is determined by the desired $(\varepsilon,\delta)$‑budget; the theorem
quantifies the exact utility penalty incurred.
\end{itemize}

\section*{Discussion}
The theorem shows that DP‑SGD retains the optimal $O(1/\sqrt{T})$ rate for
convex smooth objectives, up to an additional factor $\sqrt{d}\,\sigma$ that
captures the privacy cost. The proof follows the classic SGD analysis with
the extra stochastic term $\xi_t$ handled via its zero mean and bounded second
moment.

\newpage
\section*{Preliminaries (Definitions and Lemmas)}
\begin{lemma}[Smoothness inequality]
\label{lem:smooth}
For any $x,y\in\RR^{d}$,
\[
f(y)\le f(x)+\langle\nabla f(x),y-x\rangle+\frac{L}{2}\|y-x\|^{2}.
\]
\end{lemma}
\begin{proof}
Standard consequence of $L$‑smoothness; see e.g. Nesterov (2004). \qedhere
\end{proof}

\begin{lemma}[Bound on noisy gradient second moment]
\label{lem:noise}
For the noisy gradient $\hat g_t=\bar g_t+\xi_t$ with $\EE[\xi_t]=0$ and
$\EE[\|\xi_t\|^{2}]=d\sigma^{2}$,
\[
\EE\!\bigl[\|\hat g_t\|^{2}\mid x_t\bigr]
\le C^{2}+d\sigma^{2}.
\]
\end{lemma}
\begin{proof}
Because $\|\bar g_t\|\le C$ (average of clipped vectors) and $\xi_t$ is
independent of $x_t$,
\[
\EE\!\bigl[\|\hat g_t\|^{2}\mid x_t\bigr]
= \EE\!\bigl[\|\bar g_t\|^{2}\mid x_t\bigr]
     +\EE\!\bigl[\|\xi_t\|^{2}\bigr]
     +2\EE\!\bigl[\langle\bar g_t,\xi_t\rangle\mid x_t\bigr]
\le C^{2}+d\sigma^{2},
\]
where the cross term vanishes because $\EE[\xi_t]=0$. \qedhere
\end{proof}

\newpage
\section*{Main Proof (Step‑by‑Step Derivation)}
We prove Theorem~\ref{thm:conv}. Throughout the proof, expectations are taken
with respect to all randomness up to iteration $t$ unless otherwise noted.

\begin{proof}
\textbf{Step 1 (Smoothness expansion).}  
Using Lemma~\ref{lem:smooth} with $x=x_t$, $y=x_{t+1}=x_t-\eta\hat g_t$,
\[
f(x_{t+1})\le f(x_t)-\eta\langle\nabla f(x_t),\hat g_t\rangle
+\frac{L\eta^{2}}{2}\|\hat g_t\|^{2}. \tag{1}
\] 
[Step 1]

\textbf{Step 2 (Decompose the inner product).}  
Write $\hat g_t=\bar g_t+\xi_t$ and add/subtract $\nabla f(x_t)$:
\[
\langle\nabla f(x_t),\hat g_t\rangle
= \|\nabla f(x_t)\|^{2}
+ \langle\nabla f(x_t),\bar g_t-\nabla f(x_t)\rangle
+ \langle\nabla f(x_t),\xi_t\rangle. \tag{2}
\] 
[Step 2]

\textbf{Step 3 (Take conditional expectation).}  
Condition on $x_t$. Because $\EE[\xi_t\mid x_t]=0$,
\[
\EE\!\bigl[\langle\nabla f(x_t),\xi_t\rangle\mid x_t\bigr]=0.
\] 
Moreover, by convexity of $f$ and the definition of $\bar g_t$ (average of
clipped unbiased stochastic gradients),
\[
\EE\!\bigl[\langle\nabla f(x_t),\bar g_t-\nabla f(x_t)\rangle\mid x_t\bigr]\le 0.
\] 
Hence,
\[
\EE\!\bigl[\langle\nabla f(x_t),\hat g_t\rangle\mid x_t\bigr]
\ge \|\nabla f(x_t)\|^{2}. \tag{3}
\] 
[Step 3]

\textbf{Step 4 (Bound the noisy gradient norm).}  
From Lemma~\ref{lem:noise},
\[
\EE\!\bigl[\|\hat g_t\|^{2}\mid x_t\bigr]\le C^{2}+d\sigma^{2}. \tag{4}
\] 
[Step 4]

\textbf{Step 5 (Take expectation of (1)).}  
Using (3) and (4) in (1) and then taking full expectation,
\[
\EE[f(x_{t+1})]\le
\EE[f(x_t)]-\eta\,\EE\!\bigl[\|\nabla f(x_t)\|^{2}\bigr]
+\frac{L\eta^{2}}{2}\bigl(C^{2}+d\sigma^{2}\bigr). \tag{5}
\] 
[Step 5]

\textbf{Step 6 (Convexity inequality).}  
For convex $f$, for any optimum $x^{*}$,
\[
f(x_t)-f(x^{*})\le\langle\nabla f(x_t),x_t-x^{*}\rangle. \tag{6}
\] 
[Step 6]

\textbf{Step 7 (Relate $\|\nabla f(x_t)\|^{2}$ to distance).}  
From the update rule $x_{t+1}=x_t-\eta\hat g_t$,
\[
\|x_{t+1}-x^{*}\|^{2}
= \|x_t-x^{*}\|^{2}
-2\eta\langle\hat g_t,x_t-x^{*}\rangle
+\eta^{2}\|\hat g_t\|^{2}. \tag{7}
\] 
[Step 7]

\textbf{Step 8 (Take conditional expectation of (7)).}  
Using $\EE[\hat g_t\mid x_t]=\bar g_t$ and $\EE[\|\hat g_t\|^{2}\mid x_t]$
bounded by (4),
\[
\EE\!\bigl[\|x_{t+1}-x^{*}\|^{2}\mid x_t\bigr]
\le \|x_t-x^{*}\|^{2}
-2\eta\EE\!\bigl[\langle\bar g_t,x_t-x^{*}\rangle\mid x_t\bigr]
+\eta^{2}\bigl(C^{2}+d\sigma^{2}\bigr). \tag{8}
\] 
[Step 8]

\textbf{Step 9 (Replace $\bar g_t$ by $\nabla f(x_t)$).}  
Since $\bar g_t$ is an unbiased estimator of $\nabla f(x_t)$ after clipping,
\[
\EE\!\bigl[\langle\bar g_t,x_t-x^{*}\rangle\mid x_t\bigr]
= \langle\nabla f(x_t),x_t-x^{*}\rangle. \tag{9}
\] 
[Step 9]

\textbf{Step 10 (Combine (6) and (9)).}  
From (6) and (9),
\[
\EE\!\bigl[\langle\bar g_t,x_t-x^{*}\rangle\mid x_t\bigr]
\ge f(x_t)-f(x^{*}). \tag{10}
\] 
[Step 10]

\textbf{Step 11 (Take full expectation of (8) and use (10)).}
\[
\EE\!\bigl[\|x_{t+1}-x^{*}\|^{2}\bigr]
\le \EE\!\bigl[\|x_t-x^{*}\|^{2}\bigr]
-2\eta\,\EE\!\bigl[f(x_t)-f(x^{*})\bigr]
+\eta^{2}\bigl(C^{2}+d\sigma^{2}\bigr). \tag{11}
\] 
[Step 11]

\textbf{Step 12 (Rearrange (11)).}
\[
\EE\!\bigl[f(x_t)-f(x^{*})\bigr]
\le \frac{1}{2\eta}
\Bigl(\EE\!\bigl[\|x_t-x^{*}\|^{2}\bigr]
      -\EE\!\bigl[\|x_{t+1}-x^{*}\|^{2}\bigr]\Bigr)
+\frac{\eta}{2}\bigl(C^{2}+d\sigma^{2}\bigr). \tag{12}
\] 
[Step 12]

\textbf{Step 13 (Sum over $t=0,\dots,T-1$).}
\[
\sum_{t=0}^{T-1}\EE\!\bigl[f(x_t)-f(x^{*})\bigr]
\le \frac{1}{2\eta}\bigl(\|x_0-x^{*}\|^{2}
      -\EE\!\bigl[\|x_T-x^{*}\|^{2}\bigr]\bigr)
+\frac{T\eta}{2}\bigl(C^{2}+d\sigma^{2}\bigr). \tag{13}
\] 
[Step 13]

\textbf{Step 14 (Use non‑negativity of the norm term).}
Dropping the non‑negative $\EE[\|x_T-x^{*}\|^{2}]$ yields
\[
\sum_{t=0}^{T-1}\EE\!\bigl[f(x_t)-f(x^{*})\bigr]
\le \frac{\|x_0-x^{*}\|^{2}}{2\eta}
+\frac{T\eta}{2}\bigl(C^{2}+d\sigma^{2}\bigr). \tag{14}
\] 
[Step 14]

\textbf{Step 15 (Convert to average iterate).}  
Define $\bar x_T=\frac{1}{T}\sum_{t=0}^{T-1}x_t$.
Convexity of $f$ gives
\[
f(\bar x_T)-f(x^{*})\le\frac{1}{T}\sum_{t=0}^{T-1}\bigl[f(x_t)-f(x^{*})\bigr].
\] 
Taking expectation and applying (14):
\[
\EE\!\bigl[f(\bar x_T)-f(x^{*})\bigr]
\le \frac{\|x_0-x^{*}\|^{2}}{2\eta T}
+\frac{\eta}{2}\bigl(C^{2}+d\sigma^{2}\bigr). \tag{15}
\] 
[Step 15]

\textbf{Step 16 (Optimizing $\eta$).}  
Choosing $\eta=D/\sqrt{(C^{2}+d\sigma^{2})\,T}$ with $D\ge\|x_0-x^{*}\|$,
the RHS of (15) becomes
\[
\frac{D\sqrt{C^{2}+d\sigma^{2}}}{\sqrt{T}},
\]
which establishes the $O(1/\sqrt{T})$ rate. \qedhere
\end{proof}

\section*{Rate Derivation (With Constants)}
From (15) we have the explicit bound
\[
\boxed{
\EE\!\bigl[f(\bar x_T)-f(x^{*})\bigr]
\le
\frac{R^{2}}{2\eta T}
+\frac{\eta}{2}\bigl(C^{2}+d\sigma^{2}\bigr)},
\qquad R:=\|x_0-x^{*}\|.
\]
Setting $\eta=R/\sqrt{(C^{2}+d\sigma^{2})T}$ gives
\[
\EE\!\bigl[f(\bar x_T)-f(x^{*})\bigr]
\le
\frac{R\sqrt{C^{2}+d\sigma^{2}}}{\sqrt{T}}.
\]
Thus the constant in the $O(1/\sqrt{T})$ term is exactly $R\sqrt{C^{2}+d\sigma^{2}}$.

\section*{Special Cases}
\begin{enumerate}
\item \textbf{No privacy noise ($\sigma=0$).}  
The bound reduces to the classical SGD result:
\[
\EE[f(\bar x_T)-f(x^{*})]\le
\frac{R^{2}}{2\eta T}+\frac{\eta C^{2}}{2}.
\]
\item \textbf{High‑dimensional limit.}  
If $d$ grows while $C$ and $R$ stay fixed, the utility penalty scales as
$\sqrt{d}\,\sigma$, reflecting the well‑known “curse of dimensionality’’ for
Gaussian mechanisms.
\item \textbf{Strongly convex $f$.}  
If $f$ is $\mu$‑strongly convex, a similar derivation (using
$\|x_t-x^{*}\|^{2}\le \frac{2}{\mu}\bigl[f(x_t)-f(x^{*})\bigr]$) yields a linear
convergence term plus a steady‑state error of order $\eta(C^{2}+d\sigma^{2})$.
\end{enumerate}

\end{document}